\documentclass[11pt]{article}
\usepackage[margin=1in]{geometry}
\usepackage{amsmath,amssymb,amsthm}
\usepackage[hidelinks]{hyperref}
\setlength{\emergencystretch}{2em}

\title{The two Banach-valued ergodic questions of arXiv:0712.1740\\
\large Later affirmative answers}
\author{Autonomous mathematical discovery run \texttt{fa\_banach\_001}}
\date{17 August 2026}

\begin{document}
\maketitle

\noindent\textbf{Status.}
Literature already answered: both questions have affirmative answers, with
the standard hypotheses stated below.

\section{Source questions}

In the unnumbered ``Questions'' subsection of the outlook to
\cite[PDF pp.~19--20]{GLV}, Gruber, Lenz, and Veseli\'c ask:
\begin{enumerate}
\item whether their Banach-space-valued ergodic theorem over $\mathbb Z^d$
extends, by a similar argument, to finitely generated discrete groups of
polynomial growth; and
\item whether there is one Banach-space-valued ergodic theorem applicable to
both $\mathbb Z^d$- and $\mathbb R^d$-equivariant models, thereby avoiding a
separate treatment of discrete and continuous group actions.
\end{enumerate}

\section{Identification of the later answers}

\subsection*{Question 1: polynomial-growth groups}

Lenz, Schwarzenberger, and Veseli\'c give the requested extension in
\cite{LSV}.  Their Assumption~2.5 works on a finitely generated amenable group
with a suitable tiling F\o lner sequence, pattern frequencies, a Banach target
space, and a colouring-invariant almost-additive set function.  Their
Remark~2.6 explicitly states that every group of polynomial volume growth fits
this framework.  Theorem~3.1 then proves convergence in the Banach-space norm
and gives an explicit error estimate in terms of the boundary term and pattern
frequencies.

Thus the first question has a full affirmative answer.  The theorem covers
the requested polynomial-growth class and preserves the two characteristic
features of the source method: almost-additive Banach-valued functions and
quantitative control by tilings and frequencies.

\subsection*{Question 2: one discrete/continuous theorem}

Pogorzelski subsequently treats a second-countable, locally compact,
unimodular amenable group $G$ with its Haar measure \cite{Pog}.  The paper
notes that every discrete group and every abelian group is unimodular.
Its Theorem~5.7, ``Mean ergodic theorem for set functions,'' proves
Banach-norm convergence of bounded almost-additive functions along strong
F\o lner sequences, under a weak compactness condition on the translation
orbits and the stated compatibility with the group action.

This is precisely one abstract framework containing both cases in the second
question:
\[
 G=\mathbb Z^d \quad\text{with counting Haar measure},
 \qquad
 G=\mathbb R^d \quad\text{with Lebesgue Haar measure}.
\]
Theorem~5.7 therefore gives an affirmative answer at the level asked by the
source: the discrete/continuous distinction is absorbed into the choice of
locally compact group and Haar measure.  The later paper also supplies
pointwise and integrated-density-of-states applications, but those are not
needed for the identification.

\section{Scope and confidence}

The answer to Question~2 is an abstract unification, not a claim that every
individual model in the two earlier papers automatically satisfies every
hypothesis of Theorem~5.7.  Applications must still verify almost additivity,
measurability, the boundary condition, and weak compactness.  This is the
ordinary scope of the requested theorem rather than a gap in the answer.

The identification is high confidence.  The first later paper explicitly
singles out polynomial volume growth immediately before its Banach-valued
Theorem~3.1.  The second paper explicitly advertises Theorem~5.7 as a theorem
for locally compact unimodular amenable groups, and its preliminaries explicitly
include both discrete groups and abelian groups.  We use the corrected arXiv
version~2 of \cite{Pog}.

\begin{thebibliography}{9}
\raggedright
\bibitem{GLV}
M. J. Gruber, D. H. Lenz, and I. Veseli\'c,
\emph{Uniform existence of the integrated density of states for combinatorial
and metric graphs over $\mathbb Z^d$}, arXiv:0712.1740v1 (2007).

\bibitem{LSV}
D. Lenz, F. Schwarzenberger, and I. Veseli\'c,
\emph{A Banach space-valued ergodic theorem and the uniform approximation of
the integrated density of states}, Geom. Dedicata 150 (2011), 1--34;
arXiv:1003.3620v2.

\bibitem{Pog}
F. Pogorzelski,
\emph{Almost-additive ergodic theorems for amenable groups},
J. Funct. Anal. 265 (2013), no.~8, 1615--1666, with corrigendum/addendum;
arXiv:1211.2089v2.
\end{thebibliography}

\end{document}
